%

Our research methodology is centered around two core automated systems: an AI Scientist for generating novel scientific research and an \ourreviewer for its rigorous evaluation. These systems work in concert to explore the potential of AI in accelerating scientific discovery.

\subsection*{\ouralgo}
\ouralgo is an agentic system designed to autonomously conduct machine learning research. We present results for two modes: a template-based system that extends human-provided code, and a more open-ended template-free system that operates with much less prior guidance.
The detailed prompts used for each system are provided in Supplementary \Cref{app:v1_prompts_scientist,app:v2_prompts_scientist}. More results and analyses of the papers generated by each system are provided in \Cref{app:template_based_tables,app:case_study_v1,app:case_study_v2,appsec:template_based_papers,appsec:template_free_papers}.

\subsubsection*{Foundational Technologies}
Both versions are built upon autoregressive large language models (LLMs)~\citep{gpt4, claude3, llama3} that learn to generate text by modeling the conditional probability of a new token given preceding tokens. Through vast data and model scaling, LLMs exhibit human-like abilities, including reasoning and code generation. Agentic patterns~\citep{llm_agent_survey} such as few-shot prompting~\citep{brown2020language}
and self-reflection~\citep{shinn2024reflexion} are leveraged by \ouralgo to improve performance and reliability.
For code generation, the template-based system uses the state-of-the-art open-source coding assistant, Aider~\citep{aider}, which is designed to implement features, fix bugs, or refactor code in existing codebases.
To go further and effectively use more test-time compute, the template-free system uses LLMs to power tree search without relying on Aider.

\paragraph{Template-Based AI Scientist.}
\label{subsec:template-method}

The system is provided with a starting code template that reproduces a simple training run from a popular algorithm on a standard benchmark (e.g., training a small transformer~\citep{vaswani2017attention} on the works of Shakespeare). Its workflow unfolds in three phases:
\begin{enumerate}
    \item \textbf{Idea Generation:} The process begins with a simple experiment defined by a human-provided code template. The system then enters an iterative loop of idea generation and refinement using LLMs as a mutation operator. In each iteration, it proposes a batch of new research ideas that are variations or extensions of existing ideas in its growing archive. Each idea is a structured object containing a descriptive title, a summary of the core hypothesis, a detailed experimental plan, and self-assessed scores for interestingness (1-10 scale), novelty (1-10 scale), and feasibility (1-10 scale). This iterative growth of an idea archive is inspired by open-endedness algorithms that maintain a diverse collection of artifacts~\citep{stanley2017open, lehman2022evolutionlargemodels}. To enforce novelty, each proposed idea is automatically checked against the scientific literature via the Semantic Scholar API~\citep{fricke2018semantic}; ideas with high semantic similarity to existing work are discarded. The system is prompted to act as an ``ambitious AI PhD student who is looking to publish a paper that will contribute significantly to the field.'' For novelty assessment, the system conducts up to 10 rounds of literature search queries, with each round allowing the system to refine its search based on previous results.
    
    \item \textbf{Experiment Execution:} Once a promising idea is selected from the archive, the system devises a multi-step experimental plan with up to 5 experiments. It then executes this plan sequentially, using Aider to modify the codebase. A key feature of this phase is its robustness to runtime errors. The system automatically detects execution failures, captures the error logs, and invokes an instance of the Aider agent~\citep{aider} to perform automated debugging. The Aider agent is prompted with the failing code and the error message to generate a patch, with up to 4 re-attempt cycles per experiment. The corrected code is then used to re-run the experiment with a timeout of 7200 seconds per experiment. All experimental outcomes, including metrics, generated plots, and observations, are logged in an experimental journal. This journal serves as a form of memory, informing the subsequent steps in the experimental plan.
    
    \item \textbf{Manuscript Generation:} Upon completion of the experimental phase, the system synthesizes the findings into a full scientific paper. To do so, it uses Aider to populate 
    a standard conference \LaTeX{} template, writing sections including Introduction, Methods, Results, and Conclusion. The Results section is written by analyzing the experimental journal, summarizing key findings, and embedding the generated figures. To situate the work within the broader scientific context, the system constructs a Related Work section by querying the Semantic Scholar API for relevant literature (up to 20 search rounds) and generating summaries for each cited paper. The manuscript undergoes several passes of automated editing and refinement to improve clarity and coherence. Finally, the system compiles the \LaTeX{} source, automatically correcting any compilation errors (up to 5 correction rounds) to produce a final PDF.
\end{enumerate}

\paragraph{Template-Free AI Scientist.}
\label{subsec:template-free-method}

To overcome the limitations of a fixed starting codebase, we developed a template-free version capable of more open-ended discovery.
We use OpenAI's o3 for idea generation and code critique during experiments due to its strong reasoning capabilities, Anthropic's Claude Sonnet 4 for code generation, OpenAI's GPT-4o for cost-efficient vision-language tasks, and OpenAI's o4-mini for cost-efficient reasoning during the review stage.
This version introduces several key enhancements:
\begin{itemize}
    \item \textbf{Generalized Idea Generation:} This system's ideation process is more abstract and not tethered to an initial code implementation. It begins by generating high-level research proposals that resemble the abstract of a scientific paper. These proposals articulate a research problem, propose a novel method, and hypothesize the expected outcomes. To ensure the proposals are both grounded and novel, this process is tightly integrated with a literature review module that queries external academic databases to identify knowledge gaps and avoid rediscovering existing work. The system uses structured prompts to guide idea generation, with reflection rounds to refine proposals based on literature search results.
    (See Supplementary \Cref{app:v2_prompts_scientist} for prompts.)

    \item 
    \textbf{Experiment Progress Manager:} 
    Real-world scientific experimentation typically proceeds through distinct stages, from initial feasibility assessments to detailed ablation analyses. To emulate this structured approach, an \textbf{experiment progress manager} is introduced to coordinate four clearly defined stages of scientific experimentation:
    starting with a \textit{Preliminary Investigation} to test basic viability, moving to \textit{Hyperparameter Tuning} for optimization, then to the main \textit{Research Agenda Execution}, and concluding with \textit{Ablation Studies} to understand the contribution of different components. 
    Each stage has explicit stopping criteria. Stage 1 concludes when a basic working prototype is successfully executed. Stage 2 ends when experiments stabilize, as indicated by convergence in training curves and successful execution across at least two datasets. Stages 3 and 4 conclude when the allocated computational budget is exhausted.
    Each stage conducts its own tree search, where the specifics of this tree search process are detailed in the following bullet point.
    Each node has a maximum experiment runtime of 1 hour. At the end of each stage, an LLM-based evaluator assesses all leaf nodes and selects the most promising one to serve as the root for the next stage of exploration, effectively pruning less promising research avenues. 

    \item \textbf{Parallelized Agentic Tree Search for Experimentation:}
    To manage the complexity of open-ended research, the sequential workflow of the template-based version of \ouralgo is replaced with a parallelized agentic tree. 
    \Cref{fig:tree_ablations}A provides an overview of the approach, while \Cref{fig:tree_ablations}B shows a generated tree from an actual run.
    By default, it uses Claude Sonnet 4 for code generation. We provide results for different LLM model choices in \Cref{fig:conceptual}B.
    
Each experimental node within agentic tree search undergoes the following execution cycle: Claude Sonnet 4 first generates both a concrete experimentation plan and the associated Python code to implement the experiment. The generated code is immediately executed in a Python interpreter. If the execution encounters an error, the error message is recorded, and the node is marked as \textbf{buggy}, ending the current execution cycle for that node. If the execution succeeds, the experiment proceeds to the plotting phase. %

The system is prompted to save all relevant experimental outputs (training and validation metrics, losses, etc.) into structured numpy files. In the plotting phase, \ouralgo reads these stored results and the code, generating visualizations that summarize and illustrate the findings. These visualizations are subsequently passed to a Vision-Language Model (VLM) for critique. Any issues flagged by the VLM (such as unclear labels, missing legends, or misleading visualizations) result in the node being marked as \textbf{buggy}, and this feedback is recorded for future debugging. Nodes that successfully execute and pass the VLM review without issue are designated as \textbf{non-buggy}.

Each node is defined as a collection comprising an experiment script (e.g., a Python file), a textual description of the high-level plan implemented in the script, an execution error trace (if applicable), experiment runtime, performance metrics recorded during the experiment, code critique from OpenAI's o3 after running the script, a visualization script, file paths to the generated figures, feedback from a VLM on those figures, and the node's final status (either buggy or non-buggy).

At each iteration, the system selects several nodes from the existing tree to expand in parallel.
With a predefined probability, a \textbf{buggy node} is chosen (thus prioritizing error resolution and debugging); otherwise, a \textbf{non-buggy} node is selected for further refinement and improvement.
When choosing between non-buggy nodes, the system uses a \textbf{best-first search strategy}, guided by GPT-4o that evaluates candidates based on factors like performance metrics, training dynamics, and the quality of generated plots.
The selected nodes are expanded by creating a new child node, where the system attempts debugging if the parent node was buggy, or refines and improves upon the previous experiment if the parent was non-buggy.
Claude Sonnet 4 is used to generate the plan and experiment code for each new child node, after which all new nodes are executed concurrently in parallel, significantly accelerating the exploration process.
In addition to buggy and non-buggy nodes, the system uses specialized node variants tailored to specific experimental needs:

\begin{itemize}
\item \textbf{Hyperparameter nodes} systematically explore alternative hyperparameter configurations during Stage 2. The system maintains records of previously tested hyperparameters, preventing redundant experiments. Errors encountered during hyperparameter tuning trigger the creation of corresponding debug nodes.
\item \textbf{Ablation nodes} evaluate crucial ablation studies during Stage 4, assessing the importance of various components or assumptions underlying the experiment. Similar to hyperparameter nodes, previously tested ablation conditions are tracked to avoid repetition, and debugging nodes are created in response to any encountered errors.
\item \textbf{Replication nodes} execute replicates of their parent experiments using different random seeds. Typically, several replication nodes are created to enable the calculation of statistical measures (mean and standard deviation) of experimental outcomes, enhancing result robustness.
\item \textbf{Aggregation nodes} are special nodes created to consolidate and visualize the combined results of replication nodes. Unlike other node types, aggregation nodes do not conduct new experiments but simply generate a Python script to aggregate and summarize prior results, producing figures that explicitly show means and standard deviations.
\end{itemize}

The structured design of experiment stages and tailored node types facilitates systematic exploration across all stages.
Unlike some LLM agents that rigidly follow predefined, fine-grained workflow graphs, the \ouralgo adopts a looser structure that guides the entire empirical research cycle, enabling flexible system behavior while maintaining coherence across iterative stages.
See Supplementary \Cref{app:v2_prompts_scientist} and \Cref{app:v2_hyperparams} for the prompts and detailed hyperparameters, respectively.

    
    \item \textbf{Vision-Language Model (VLM) Integration:} This system incorporates Vision-Language Models (VLMs) using GPT-4o 
    to analyze and provide feedback on visual data. During experimentation, generated plots are fed to a VLM, which is prompted to act as a scientist and critique them. 
    For example, it might flag nonsensical axes, issues in the quality of generated examples, or suggest clearer ways to present the data.
    This feedback is used to generate new experimental nodes in the tree search aimed at addressing the identified issues. During manuscript preparation, the VLM assesses the alignment between figures and their corresponding captions, ensuring that the caption accurately describes the plot and highlights the key takeaways, thus improving the overall quality and clarity of the paper. The VLM reviews include detailed analyses of figure content, caption accuracy, and integration with the main text.
    (See Supplementary \Cref{app:v2_prompts_scientist} for prompts.)

    \item \textbf{Generalized Dataset Access:} 
    To broaden its research capabilities, the system is prompted to dynamically integrate datasets from public repositories by formulating queries to the HuggingFace Hub~\citep{wolf2020huggingface}. 
    A set of 10 example datasets available on HuggingFace is listed in the prompt, and the system can automatically generate the necessary data-loading code to use a selected dataset in its experiments. 
    This approach partially relaxes the constraint of working with a fixed, pre-defined set of datasets by allowing human scientists to easily update the candidate list.
    For datasets not available on HuggingFace, human scientists can download them from public data repositories (e.g., open-access archives), store them locally, and add usage instructions to the prompt. These locally stored datasets can then be used alongside HuggingFace datasets by \ouralgo.
    (See Supplementary \Cref{app:v2_prompts_scientist} for prompts.)
    
    \item \textbf{Enhanced Manuscript Writing:} The template-free system moves away from the incremental Aider-based approach to direct LaTeX generation using a reasoning model such as OpenAI's o1~\citep{OpenAIOS} followed by reflection~\citep{shinn2023reflexion}. The system first aggregates experimental results from multiple stages into compound figures using a dedicated plot aggregation step. The manuscript writing process includes specific prompts for different workshop formats (e.g., the ICBINB workshop focusing on negative results), with detailed guidelines for each section, including the title, abstract, introduction, methods, experiments, and conclusions. The system undergoes multiple reflection cycles, incorporating feedback from LaTeX linters and VLM reviews to improve figure quality and text-figure alignment. (See our code and Supplementary \Cref{app:v2_prompts_scientist} for prompts and full details.)
\end{itemize}

The complete generation process for the template-free system typically takes several hours to over 15 hours, depending on problem complexity.


\subsection*{\ourreviewer}
To assess the quality of the AI-generated research, we built an \ourreviewer using o4-mini~\citep{openai2025o3o4mini}. This component is designed to emulate the peer-review process of a top-tier machine learning conference by adhering to the official NeurIPS reviewer guidelines. The agent processes the manuscript's PDF to produce a structured review, including numerical scores for soundness, presentation, and contribution, along with a list of strengths and weaknesses and a preliminary accept/reject decision (Supplementary \Cref{app:reviewer_details}). All prompts used for \ourreviewer are provided in Supplementary \Cref{app:v1_prompts_reviewer}. 

\paragraph{Review Process.}
The \ourreviewer follows a multi-stage process. First, the system is prompted with the role: ``You are an AI researcher who is reviewing a paper that was submitted to a prestigious ML venue.'' The review prompt provides the paper content along with detailed NeurIPS reviewer guidelines and asks for a structured JSON response including summary, strengths, weaknesses, questions, limitations, ethical concerns, and numerical scores (soundness, presentation, contribution, overall score 1-10, and confidence level). To improve robustness, the final assessment is a meta-review that ensembles five independent reviews. The five reviews are generated for each paper and aggregated into a single meta-review, with an LLM taking the role of an Area Chair to find consensus among the individual reviews.

\paragraph{Validation.}

We benchmarked the \ourreviewer against human decisions using ICLR data from the publicly available OpenReview dataset~\citep{gonzalez2024learning}. 
The \ourreviewer achieves a comparable balanced accuracy with humans (69\% versus 66\%; see Supplementary \Cref{app:v1_reviewer_details} for details) and a higher F1-score compared to inter-human group agreement (0.62 versus 0.49) in the NeurIPS 2021 consistency experiment~\citep{beygelzimer2021neurips}, where roughly 10\% of submissions were randomly selected and sent to two independent review committees, providing a real-world benchmark of inter-reviewer consistency (\Cref{tab:reviewers_main_paper}).
These results suggest that LLM-based agents can provide valuable feedback that aligns with the average human expert opinion.
We highlight that there was a different set of paper submissions in the ICLR and NeurIPS paper pools and thus distribution shift, so this comparison is not exact.
However, ICLR is the only major ML conference that releases all accept and reject decisions for us to perform the analysis, and the NeurIPS 2021 experiment is the only modern version of the human consistency experiment, and is thus the only possible comparison.

\newpage
\section*{Data Availability}

For the nanoGPT~\citep{karpathy2022nanogpt} experiments, the template-based version of \ouralgo used the Shakespeare character~\citep{char_shakespeare}, enwiki8~\citep{hutter_prize} and text8~\citep{text8} datasets.
The template-free version of \ouralgo used the Crop Pest and Disease Detection dataset~\citep{mensah2023crop} for one of the papers submitted to the ICLR workshop experiment, and the Waterbirds~\citep{waterbirds_dataset} and CelebA~\citep{liu2015faceattributes} datasets for the experiments shown in \Cref{fig:conceptual}B, \Cref{fig:tree_ablations}B, and \Cref{fig:tree_ablations}C.
In all other cases, the template-free version used datasets available through the HuggingFace Hub~\citep{wolf2020huggingface}.


\section*{Code Availability}
\label{sec:code_availability}

The code for the template-based version of \ouralgo and \ourreviewer is available at \url{https://github.com/SakanaAI/AI-Scientist}.
Additionally, the code for the template-free version of \ouralgo is available at \url{https://github.com/SakanaAI/AI-Scientist-v2}.
Both code repositories are licensed under the Apache License 2.0.


\section*{Acknowledgements}
The authors would like to thank Irene Zhang, Johannes von Oswald, Takuya Akiba, Yujin Tang, Kosuke Nakago, Kou Misaki, Haruki Goda, Yuichi Inoue, Aaron Dharna, Ben Norman, Jenny Zhang, Anna Olerinyova, and Felicitas Muecke-Wegner for helpful feedback.
This work was supported by grants from Schmidt Futures, NSERC, the Vector Institute, the Canada CIFAR AI Chairs program, and a donation from Rafael Cosman.

\section*{Author Contributions}

There are four authors who contributed equally to this work, and are listed in alphabetical order.

\textbf{Yutaro Yamada [shared first, equal corresponding]:} Co-led the project and contributed core ideas. Coded the core tree-search and template-free version of \ouralgo. Ran paper generation experiments. Read and validated the work of many AI-generated papers to select submissions and checked the paper code implementations. Led the writing of the paper. Wrote detailed analyses of the submitted papers for our manuscript.

\textbf{Robert T. Lange [shared first, equal corresponding]:} Co-initiated, co-led the project, and contributed core ideas. Conceived and coded core parts of the \ourreviewer, tailored the paper generation pipeline to the workshop, and ran the paper generation experiments. Organized the workshop communication process. Read and validated the work of many AI-generated papers to select submissions and checked the paper code implementations. Led the writing of the paper. Wrote detailed analyses of the submitted papers for our manuscript.

\textbf{Cong Lu [shared first, equal corresponding]:} Initiated the project, co-led the project, and conceived the original ideas for \ouralgo, including the use of SWE agents like Aider to execute scientific ideas autonomously. Coded core parts of the idea generation, \ourreviewer, tool use, experiment aggregation, and paper writing framework. Wrote and led the IRB approval process for the workshop experiments, and evaluated AI-generated paper submissions. Led the writing of the paper.

\textbf{Chris Lu [shared first, equal corresponding]:} Initiated the project, co-led the project. Conceived the original idea and structure for \ouralgo and developed the first working system, which demonstrated autonomous end-to-end paper generation. Conceived the experiment setup and evaluations in the original preprint. Led the development of the template-based version of \ouralgo. Ran paper generation experiments. Led the writing of the paper. Provided advice, feedback, and writing.

\textbf{Shengran Hu:} Enhanced the iterative \ourreviewer with VLM feedback, contributed to the experiment and paper writing framework, ran \ourreviewer benchmark and ablation experiments, helped run paper generation experiments, read and validated the work of many AI-generated papers to select submissions, and checked the paper code implementations. Helped write and iterate over drafts of the paper. Helped write the IRB approval.

\textbf{Jakob Foerster:} Provided advice, feedback, and writing.

\textbf{David Ha [equal corresponding]:} Provided overarching guidance for the research project, offering technical insight, advice, feedback, and writing. Oversaw the public communication process.

\textbf{Jeff Clune [equal corresponding]:} Provided overarching guidance for the research project, offering technical insight, advice, feedback, and writing. Oversaw the IRB application process. Evaluated AI-generated paper submissions.

\section*{Competing Interests}
The Principal Investigator, Jeff Clune, has affiliations with the Vector Institute and Google DeepMind.
This project has Vector Institute affiliations but is not connected to Google DeepMind.
Several co-authors are employees of or consultants for Sakana AI, a machine learning research company involved in the design of \ouralgo.
The Principal Investigator is not financially compensated by Sakana AI.
These arrangements have been reviewed and approved by the University of British Columbia in accordance with its conflict of interest policies.

\section*{Ethics Approval}
This study received ethics approval from The University of British Columbia Behavioral Research Ethics Board (BREB) (protocol number H24-02652).
The research was conducted in full cooperation with the ICLR conference leadership and the relevant workshop organizers.
In accordance with the approved protocol, human participants (peer reviewers) were informed that a small number of submissions to the workshop were AI-generated, although not which specific papers.
Participants had the option to opt out of reviewing any potentially AI-generated manuscripts.
All AI-generated submissions were withdrawn following the review process, regardless of the outcome.

